\documentclass[11pt]{article}
\usepackage[margin=1in]{geometry}
\usepackage{amsmath,amssymb,amsthm,booktabs,mathtools,microtype}
\usepackage{xurl}
\usepackage[colorlinks=true,linkcolor=blue,urlcolor=blue]{hyperref}
\setlength{\emergencystretch}{2em}
\newtheorem{theorem}{Theorem}
\newtheorem{lemma}[theorem]{Lemma}
\newtheorem{corollary}[theorem]{Corollary}
\newcommand{\dd}{\,\mathrm d}
\newcommand{\Q}{\mathbb Q}
\title{Completing Atlas 130 and Atlas 203\\
\large Exact rational certificates for both missing density intervals}
\author{}
\date{September 21, 2026}
\begin{document}
\maketitle

\begin{abstract}
We prove the chromatic-polynomial graphon lower bound for the two remaining
partial rows of the six-vertex catalogue.  New rational certificates cover
$[1/2,2/3]$ for Atlas 130 and $[2/3,3/4]$ for Atlas 203.  Together with the
previously verified high-density certificates, these prove both inequalities
on their full required intervals, giving 94 positive and 23 negative rows.
Each new certificate is an exact identity involving integrals of rooted
squares, fixed-density identities, and nonnegative induced-density terms.
A separate verifier reconstructs all counting coefficients using integer
enumeration, proves positivity by exact rational Schur complements, and checks
936 rational equations per certificate.  No numerical solver output is used
as a premise of the proof.  Both full bounds now have complete Lean proofs. Fresh sequential builds
use 100 row-specific files each and meet the one-hour, 16 GiB limits;
the shared infrastructure is measured separately. The complete theorem
axiom audits admit only the three standard Lean axioms. See
\path{lean/docs/ATLAS130_VERIFICATION_PROGRESS.md} and
\path{lean/docs/ATLAS203_VERIFICATION_PROGRESS.md} for the records.
\end{abstract}

\section{The complete results}

Let $W\colon\Omega^2\to[0,1]$ be a symmetric measurable kernel on a
probability space, and put $p=\int_{\Omega^2}W\dd\mu^2$.  All densities are
homomorphism densities, with the target at $p=1$ interpreted polynomially.

\begin{theorem}\label{thm:complete}
For every such graphon,
\begin{align}
 t(H_{130},W)&\ge p^3(2p-1)^2 &&(1/2\le p\le1),\label{eq:130}\\
 t(H_{203},W)&\ge p(2p-1)(3p-2)(10p^2-14p+5)
       &&(2/3\le p\le1).\label{eq:203}
\end{align}
\end{theorem}

Atlas 130, with graph6 string \verb|E{CW|, consists of the triangles $012$
and $345$ joined by the edge $03$.  Its edges and chromatic polynomial are
\[
 E_{130}=\{01,02,03,12,34,35,45\},\qquad
 P_{130}(z)=z(z-1)^3(z-2)^2.
\]
Atlas 203, with graph6 string \verb|E~^G|, has
\[
 E_{203}=\{01,02,03,05,12,13,14,15,23,24,34,45\}.
\]
Equivalently it is $K_6$ with $04,25,35$ removed.  Its chromatic polynomial is
\[
 P_{203}(z)=z(z-1)(z-2)(z-3)(z^2-6z+10).
\]
Substitution of $z=1/(1-p)$ gives exactly the targets in
\eqref{eq:130}--\eqref{eq:203}.  The independent checker also recovers each
target by enumerating all independent-set partitions of the six source
vertices, using $P_H(z)=\sum_\pi(z)_{|\pi|}$.

\section{An arbitrary-graphon certificate lemma}

Here we give all analytic ingredients of the computer-assisted proof.
There is no extremal-graphon assumption or finite-host approximation.

\subsection{Root patterns and rooted probabilities}

For $m\in\{0,2,4\}$, set $r=(6-m)/2$.  Take one labelled representative
$\sigma$ of each unlabelled graph on $m$ vertices, using the NetworkX Atlas
order.  There are respectively $1,2,11$ representatives, hence 14 root types.
For $x=(x_0,\ldots,x_{m-1})$, write
\[
 w_\sigma(x)=\prod_{ij\in E(\sigma)}W(x_i,x_j)
 \prod_{\substack{i<j<m\\ij\notin E(\sigma)}}(1-W(x_i,x_j)).
\]
For the empty type this weight is 1.

Conditional on the root points, sample $r$ further independent points and
independent Bernoulli edge indicators for all pairs having at least one
new endpoint.  Each pair $uv$ has success probability $W(x_u,x_v)$.
Group the resulting edge patterns into classes under permutations of the
$r$ new vertices, fixing the roots individually.  Let $f_\sigma(x)$ be the
vector of the probabilities of these classes.  Thus its entries sum to 1;
they do not include the root weight $w_\sigma$.  The vector sizes are
\[
 \begin{array}{c|ccc}m&0&2&4\\ \hline
 r&3&2&1\\ \dim f_\sigma&4&20&16.\end{array}
\]
For a precise encoding, order the pairs $i<j<m+r$ with $j\ge m$
lexicographically, encode an edge set by its binary mask in that order,
replace it by the least mask under permutations of the new vertices, and
order the distinct least masks increasingly.

For a rational positive-semidefinite matrix $Q$,
\begin{equation}
 \int_{\Omega^m}w_\sigma(x)f_\sigma(x)^{\mathsf T}
 Qf_\sigma(x)\dd\mu^m(x)\ge0.\label{eq:square}
\end{equation}
Indeed, the integrand is a nonnegative weight times a nonnegative quadratic
form.  This also holds with polynomially varying positive-semidefinite $Q$.

\subsection{The independent six-point expansion}

Let $\mathcal G_6$ be the 156 unlabelled simple graphs on six vertices, and
let $\rho_J(W)$ be the probability that six independent points, with the
conditionally independent edge indicators just described, induce a graph
isomorphic to $J$.  In particular, $\rho_J(W)\ge0$ and
$\sum_J\rho_J(W)=1$.

For a graph $F$ on at most six vertices, let $h_F(J)$ be the proportion of
injective maps $V(F)\to V(J)$ that preserve every edge of $F$.  Exchangeability
and conditional independence give the exact identity
\begin{equation}
 t(F,W)=\sum_{J\in\mathcal G_6}\rho_J(W)h_F(J).
 \label{eq:hom-expansion}
\end{equation}
Injectivity here concerns the six sampled \emph{indices}, not distinctness
of their underlying points in $\Omega$.  Thus this is exact even on an atomic
probability space, and does not replace graphon homomorphism densities by
finite-host injective densities.

Similarly let $M^\sigma_{ij}(J)$ be the following unconditional probability.
Choose a uniformly random ordering of the six vertices of $J$.  The first
$m$ vertices must induce the specific labelled type $\sigma$; the next $r$
and final $r$ vertices must give rooted classes $i$ and $j$, respectively.
Unspecified edges between these two branch sets are ignored.  Then
\begin{equation}
 \int w_\sigma f_{\sigma,i}f_{\sigma,j}\dd\mu^m
   =\sum_{J\in\mathcal G_6}\rho_J(W)M^\sigma_{ij}(J).
 \label{eq:gram-expansion}
\end{equation}
To prove this, first condition on the roots.  The two branch experiments
are independent and the root edges are sampled once.  Averaging the resulting
product over all sampled six-point graphs proves the identity.  In
particular, $M^\sigma$ is \emph{not} divided by the number of copies of
$\sigma$ in $J$; that conditional normalization would invalidate this proof.

\subsection{The interval identity}

For $a<b$, put $p(s)=a+(b-a)s$ and $0\le s\le1$.  Write
\[
 B_{d,k}(s)=\binom dk s^k(1-s)^{d-k}.
\]
For each root type, the certificate contains integer matrices $L_\sigma$,
$R_\sigma$, $C_{0,\sigma}$, $C_{1,\sigma}$ and symmetric rational
\emph{positive-definite} matrices $Z_{0,\sigma}$, $Z_{1,\sigma}$.
Define
\begin{align*}
 A_\sigma(s)&=[(1-s)L_\sigma\ \ sR_\sigma]C_{0,\sigma},\\
 Q_\sigma(s)&=A_\sigma(s)Z_{0,\sigma}A_\sigma(s)^{\mathsf T}
   +s(1-s)C_{1,\sigma}Z_{1,\sigma}C_{1,\sigma}^{\mathsf T}.
\end{align*}
Every $Q_\sigma(s)$ is positive semidefinite on the closed interval.
The certificate also contains arbitrary rational coefficients $u_{G,k}$
for each of the 11 four-vertex graphs $G$, and nonnegative rationals
$z_{J,k}$ for $J\in\mathcal G_6$.  Put
\[
 \gamma_G(s)=\sum_{k=0}^4u_{G,k}B_{4,k}(s),\qquad
 z_J(s)=\sum_{k=0}^5z_{J,k}B_{5,k}(s)\ge0.
\]
The exact identity checked by the verifier is
\begin{align}
 t(H,W)-\Phi_H(p(s))
 &=\sum_\sigma\int w_\sigma f_\sigma^{\mathsf T}
       Q_\sigma(s)f_\sigma\dd\mu^m\notag\\
 &\quad+\sum_{G\in\mathcal G_4}\gamma_G(s)
       \bigl(t(G\sqcup K_2,W)-p(s)t(G,W)\bigr)\notag\\
 &\quad+\sum_{J\in\mathcal G_6}z_J(s)\rho_J(W).
 \label{eq:certificate}
\end{align}
This identity holds for every graphon and every real $s$ as a polynomial
identity; the positivity conclusion uses $s\in[0,1]$ and $t(K_2,W)=p(s)$.

\begin{lemma}\label{lem:certificate}
If the stated rational positivity conditions and identity
\eqref{eq:certificate} hold, then $t(H,W)\ge\Phi_H(t(K_2,W))$ whenever
$a\le t(K_2,W)\le b$.
\end{lemma}
\begin{proof}
Choose $s=(t(K_2,W)-a)/(b-a)$.  All square terms and all induced-density
terms are nonnegative.  Independence on the disjoint components gives
$t(G\sqcup K_2,W)=t(K_2,W)t(G,W)$, so every middle term is zero, regardless
of the signs of $\gamma_G$.  This proves the result, including both endpoints.
\end{proof}

\section{Exact verification and interval coverage}

The two exact data files, relative to \path{experiments/}, are
\begin{center}
\footnotesize
\begin{tabular}{@{}cclrr@{}}
\toprule
Atlas & interval & filename & zero slacks & positive slacks\\
\midrule
130 & $[1/2,2/3]$ & \path{atlas130_exact_lower_induced_interval.json} &43&893\\
203 & $[2/3,3/4]$ & \path{atlas203_exact_lower_induced_interval.json} &31&905\\
\bottomrule
\end{tabular}
\end{center}
Each file contains the matrices and all coefficients in
\eqref{eq:certificate}, as integers or exact rational strings.  There are
28 positive-definite rational Gram factors per certificate.

The standalone checker is
\path{experiments/verify_remaining_cases_induced.py}.  It performs the
following exact checks:
\begin{enumerate}
\item Rebuild the Atlas source graph, all 156 hosts, the 14 root types, and
the rooted class encodings.  Recover the target by independent-set partitions.
\item Verify symmetry and positivity of all 28 Gram factors by rational
Schur elimination: every pivot is strictly positive.
\item Rebuild $h_F(J)$ and $M^\sigma_{ij}(J)$ by enumerating all $6!=720$
ordered samples in each host.  All coefficients are integer counts over 720.
\item Expand both sides of \eqref{eq:certificate} in the degree-five
Bernstein basis.  Verify all $156\cdot6=936$ equations over $\Q$, and
verify that all 936 residual coefficients are nonnegative.
\end{enumerate}
Equations \eqref{eq:hom-expansion} and \eqref{eq:gram-expansion} show why these
finite checks prove the entire graphon identity.  This is not a test of the
inequality on 156 particular graphons.

From the \path{exhaustive_research/} directory, run
\begin{verbatim}
python experiments/verify_remaining_cases_induced.py \
  experiments/atlas130_exact_lower_induced_interval.json \
  experiments/atlas203_exact_lower_induced_interval.json
\end{verbatim}
On PowerShell this can be entered on one line.  The verifier uses NetworkX
and python-flint, together with Python's exact integer and fraction arithmetic.
It imports no discovery scripts, solver, or numerical cache.

\begin{proof}[Proof of Theorem~\ref{thm:complete}]
The two exact checks establish the hypotheses of
Lemma~\ref{lem:certificate}, proving \eqref{eq:130} on $[1/2,2/3]$ and
\eqref{eq:203} on $[2/3,3/4]$.
The earlier four-root certificates, described in
\path{notes/s4_exact_interval_sos_remaining_cases.tex}, prove the respective
inequalities on $[2/3,1]$ and $[3/4,1]$.  Their exact data are
\path{experiments/atlas130_exact_upper_interval_sos.json} and
\path{experiments/atlas203_exact_upper34_interval_sos.json}; both were
rechecked with \path{experiments/s4_interval_verify.py} in this completion.
The closed intervals meet and cover the full required ranges.
\end{proof}

\section{What changed relative to the earlier attempts}

The new search allows every induced root pattern on zero, two, or four
roots, retains a nonnegative induced six-point residual, and includes the
exact fixed-density identities with polynomial multipliers.  The earlier
restricted four-root search did not supply these particular lower-interval
certificates.  We make no claim that its cone cannot do so.

For Atlas 130, this proof does not assert positive correlation of the two
rooted triangles, discard complement-codegree slack, or prove the stronger
conjecture $t(H_{130},W)\ge(2p-1)t(F,W)$ from the earlier operator note.
That stronger conjecture remains unproved here.  The invalid routes in the
failed-attempts ledger therefore stay invalid even though the target is now
settled.  For Atlas 203, the proof does not infer positivity of a cone from
positivity of its base; indeed its base Atlas 48 is negative.

For reproducibility, the discovery and reconstruction scripts are
\path{experiments/remaining_cases_flag_search.py},
\path{experiments/remaining_cases_interval_search.py}, and
\path{experiments/remaining_cases_rationalize.py}.  Floating-point solves
suggested rational subspaces; rounding followed by exact linear corrections
produced the saved data.  No inferred kernel or numerical tolerance needs
to be trusted: the independent checker tests the resulting rational factors
and the complete identity directly.

\end{document}
